一层密度为 \(\rho>0\) 的不可压缩、无黏、无旋理想流体占据 \(-h<z<\eta(x,t)\)，底部为水平不可渗透固壁 \(z=-h\)，\(h>0\)。重力为 \(g>0\) 向 \(-z\)。只研究小幅度自由表面扰动 \(\eta(x,t)\)，忽略表面张力和流动非线性。速度场可写成 \(\boldsymbol v=\nabla\Phi\)，取沿 \(x\) 传播的单色波因子 \(e^{i(kx-\omega t)}\)，\(k>0\)。
(1) 由不可压与底部不透流得到速度势满足的空间方程与底部边界条件；(2) 从线性化的自由表面运动学和动力学条件推导精确色散关系 \(\omega(k)\)；(3) 求相速度 \(c_p\) 与群速度 \(c_g\)，证明二者比值的解析式；(4) 给出 \(kh\ll1\)、\(kh\gg1\) 两个极限的色散关系、速度和传播含义，说明在浅水极限为何波包近似无频散。
